\subsection{Mixed event on MuDst}
\label{sec:mudst}

The picoDst stores a time only for BTOF hits that production matched to a track (\mtTimedFracVeight\% of the hits,
Sec.~\ref{sec:bkg}). The unmatched hits therefore cannot compete in a picoDst mixed event, and the late same-event
background could not be reproduced (best picoDst result: SE/ME = \mtSbRatioAllVnine{} between the p and d peaks).
The MuDst keeps every \fcmd{StMuBTofHit} with its leading and trailing edge, so this section repeats the mixed event on
MuDst. Two methods are used: an emulation of the production matching in our own maker (Sec.~\ref{sec:muemul}), and, as the
main result, a re-run of the production makers themselves on the overlaid events (Sec.~\ref{sec:muov}).

\paragraph{Sample.} The P24iy MuDst are on HPSS only. Six MuDst files of three runs were restored through the Data
Carousel (\fpath{scripts/mudst_restore.sh}, list \fpath{filelists/mudst/restore_test.txt}); they are the raw files of the
picoDst sample. The two files of run 20088015 are not used: reading them floods the log with trigger-data decoding errors
and the jobs stall. The four files used are \fpath{st_physics_20060065_raw_1000011}, \fpath{_5500010},
\fpath{st_physics_20073019_raw_4500013} and \fpath{_7000013}: \ovtEvAll{} events, of which \ovtEvMtwo{} pass the
event selection of Sec.~\ref{sec:m2} and have a VPD start time. The m$^2$ numbers in this section use these events only;
they are a test sample of two runs. Statistical uncertainties only.

\subsubsection{Times of matched and unmatched hits}
\label{sec:muprobe}
\fpath{macros/probeMuBTof.C} reads all hits of the four files (\prEvents{} selected events) in two passes. Pass 1 measures a
per-cell offset as the median of $t_{\rm LE}-t_{\rm tof}-T_0$ over hits matched by production, where $t_{\rm tof}$ is the
production time of flight and $T_0$ the event median of $t_{\rm LE}-t_{\rm tof}$. Pass 2 gives every hit the time
$t = t_{\rm LE} - {\rm off}_{\rm cell} - T_0$. On matched hits $t$ reproduces the production time of flight with an RMS of
\prCalRMS\,ns (\prCalOneNs\% within 1\,ns); \prNoOffFrac\% of the unmatched hits are in cells without an offset and are
dropped. Of the hits, \prNMatched{} are matched and \prNUnmatched{} (\prUnmatchedFrac\%) are unmatched.

Unmatched hits are systematically later (Fig.~\ref{fig:probe_time}). The fraction at 8--10\,ns is \prMpeak\% for matched
and \prUpeak\% for unmatched hits; at 15--26\,ns it is \prMmid\% and \prUmid\%, and above 26\,ns \prMlate\% and \prUlate\%.
Per ns, unmatched hits are \prRatioLate{} times more frequent than matched hits at 30--60\,ns. The late fraction of
unmatched hits falls only mildly with the BTOF hit multiplicity (\prLateLoMult\% for the lowest, \prLateHiMult\% for
the highest classes). Their ToT distribution has the same mean as for matched hits (\prTotU{} vs \prTotM\,ns). This is the
behaviour of slow particles that are not reconstructed in the TPC, not of electronic noise; the hits are in the same bunch
crossing by construction of the production.

All matched hits have $10<{\rm ToT}<30$\,ns (\prTotMout\% outside, Fig.~\ref{fig:probe_time} right), while \prTotUlow\% of
the unmatched hits are below 10\,ns and \prTotUhigh\% above 30\,ns. Hits outside this window cannot end up as a calibrated
match, so they never contribute to random matches.

\begin{figure}[H]
\centering
\figpanel{0.98}{probe_time}{BTOF hits of the four MuDst files, normalized to unit area. Left: calibrated hit time $t$
(0.5\,ns bins, log scale). Right: ToT (0.2\,ns bins). Blue: hits matched to a track by production; red: unmatched hits.
The legend of the left panel applies to both. Unmatched hits have a long late tail: \prUlate\% above 26\,ns against
\prMlate\% for matched hits.}
\caption{Time and ToT of BTOF hits in MuDst (\fpath{scripts/make_probe.C}).}
\end{figure}

\subsubsection{Closure test: production matching and calibration re-run on MuDst}
\label{sec:muclos}
The SL24y makers \fcmd{StBTofMatchMaker}, \fcmd{StVpdCalibMaker} and \fcmd{StBTofCalibMaker} can run on MuDst
(\fcmd{setMuDstIn()}). The chain (\fpath{macros/runTofClosure.C}) is \fcmd{StMuDstMaker}, \fcmd{St_db_Maker} (current
database), the three makers with the y2019 geometry, and \fpath{StRoot/StTofClosureMaker}. That maker stores the
production matching of every global track before the re-run and compares it afterwards.
On \clEvents{} events of \fpath{st_physics_20060065_raw_1000011}:
\begin{itemize}
\item all \clNSame{} production matches are found again in the same cell; there are \clNOther{} matches in another cell,
  \clNOnlyProd{} only in production and \clNOnlyRerun{} only in the re-run;
\item the re-run time of flight differs from production by a median of \clDtofMed\,ps (68\% between \clDtofLo{} and
  \clDtofHi\,ps, 98\% within $\pm$\clDtofNinetyEight\,ps); \clDbetaCore\% of the tracks agree in $1/\beta$ within 0.005;
\item the event start time differs by a median of \clDtsMed\,ps, and \clDtsSame\% of the events agree within 5\,ps.
\end{itemize}
The residual differences are small compared with the BTOF resolution; they are attributed to calibration tables updated
after the production (not verified: the database time stamp was not pinned to the production).

\choice{The re-run of the production BTOF makers on MuDst reproduces the production matching exactly and the time of flight
to within tens of ps. It is therefore used for the mixed event in Sec.~\ref{sec:muov}, so that same event and mixed event
share the same matching and calibration.}

\subsubsection{Method 1: emulation of the production matching}
\label{sec:muemul}
\fpath{StRoot/StTofMEMuMaker} (derived from \fpath{StTofMEMaker}, same projection and mixing classes) emulates the SL24y
matching rule on the overlaid hit map. Every production-valid global track (flag, $n_{\rm HitsFit}\ge15$, ratio $\ge0.52$,
outer helix) is projected; candidates are the projected cell and its two neighbours in the module. A cell claimed by more
than one track is discarded, and among several candidates the largest ToT (below 40\,ns) wins. Hit times are our own
calibration ($t_{\rm LE}$ minus per-cell offset minus event $T_0$, RMS \muCalRMS\,ns against production). On the same event,
the emulation picks the production hit for \muChoiceSame\% of the production matches.

In the mixed event, the hits of a partner event are overlaid on the hit map of the event, and hits of the two events in
the same cell are merged (earlier leading edge, larger ToT). The partner's own production-valid tracks also claim the
cells of the partner hits. A partner hit made by a reconstructed particle is thus discarded as shared, exactly as in its own
event. Without these claims the random fraction is \muFrTimeCi\% in 0--5\% (green line in Fig.~\ref{fig:mu_m2_cent}).
With them it is \muFrProdCi\% (0--5\%) and \muFrProdAll\% (0--80\%). The same-event histogram uses the production $\beta$.

\begin{figure}[H]
\centering
\figpanel{0.9}{mu_m2_cent}{$m^2$ of QA tracks, 0--5\%, nine momentum bins, four MuDst files (log scale). Black points:
same event (production $\beta$). Red filled: mixed event with the production emulation including the partner-track claims,
normalized per QA track. Green line: the same without partner-track claims. Between the p and d peaks SE/ME =
\muHiProdCi{} (0.2--3\,GeV/$c$).}
\caption{MuDst mixed event, emulation of the production matching, 0--5\% (\fpath{scripts/make_mu.C}).}
\end{figure}

\caveat{The emulation uses our own hit times (RMS \muCalRMS\,ns, no slewing or $z_{\rm local}$ correction). Random matches to
pion-like partner hits are smeared about three times wider in $m^2$ than in the same event. In the thin bins just below and
beside the pion and kaon peaks the mixed event exceeds the same event bin by bin, up to ME/SE = \muVmaxProdCi{} in 0--5\%
($0.4<p<1.5$\,GeV/$c$). This is the motivation for method 2.}

\subsubsection{Method 2 (main result): production makers re-run on overlaid events}
\label{sec:muov}
\fpath{StRoot/StTofMEOverlayMaker} (macro \fpath{macros/runTofMEOverlay.C}, one HTCondor job per file and setup) runs after
the chain of Sec.~\ref{sec:muclos}, which first re-runs the three production makers on the event itself; this gives the same
event. Mixing classes are 14 $V_z$ bins $\times$ 9 refMult classes per run, with the last 5 events of a class as partners.
For each partner:
\begin{enumerate}
\item the partner hits are appended to the \fcmd{StMuBTofHit} array of the event, with leading and trailing edges shifted
  by $t_{\rm start}^{\rm event}-t_{\rm start}^{\rm partner}$ (VPD start times from \fcmd{StBTofHeader}), so that each
  partner hit keeps its time of flight. A partner hit in a cell that already has an overlapping hit
  ($|\Delta t_{\rm LE}|$ below the ToT of the earlier hit) is merged into it: earlier leading edge, later trailing edge;
\item the partner's global tracks with $n_{\rm HitsFit}\ge15$ and $0<{\rm flag}<1000$ (on average \ovtNTrkAdd{} per partner)
  are appended to the global-track array. They have no primary-track link and are not counted; their role is to claim the
  cells of the partner hits, so that a partner hit made by a reconstructed particle is discarded as shared, as in its own
  event;
\item \fcmd{StBTofMatchMaker}, \fcmd{StVpdCalibMaker} and \fcmd{StBTofCalibMaker} are run again on the overlaid event.
  The match maker resets the BTOF information of every track and hit at the start of each pass;
\item a QA track of the event matched to an appended hit, or to a merged hit whose leading edge came from the partner, is a
  random match. Its $m^2$ is computed from the re-run time of flight and path length, exactly as for the same event;
\item the appended hits and tracks are removed and the hit array of the event is restored.
\end{enumerate}
The mixed event is normalized per QA track in each (centrality, $p$) bin. The same event and the mixed event therefore share
the production matching rule (shared cells, largest ToT), the production calibration and the same path-length formula.

\begin{figure}[H]
\centering
\figpanel{0.9}{ovtrk_m2_cent}{$m^2$ of QA tracks, 0--5\%, nine momentum bins, four MuDst files (log scale). Black points:
same event. Red filled: mixed event with the production makers re-run on the overlay of partner hits and partner tracks,
normalized per QA track. Both use the re-run calibration. Random matches are \ovtFrCi\% of the TOF matches; SE/ME between
the p and d peaks is \ovtHiCi{} (0.2--3\,GeV/$c$) and \ovtHiMinCi--\ovtHiMaxCi{} per bin at 0.4--2\,GeV/$c$.}
\caption{MuDst mixed event (main result), production makers re-run on overlaid hits and tracks, 0--5\% (\fpath{scripts/make_ov.C}).}
\end{figure}

\begin{figure}[H]
\centering
\figpanel{0.9}{ovtrk_m2_all}{As Fig.~\ref{fig:ovtrk_m2_cent} for 0--80\% centrality. Random matches \ovtFrAll\%; SE/ME between
the p and d peaks \ovtHiAll{} (0.2--3\,GeV/$c$) and \ovtHiMinAll--\ovtHiMaxAll{} per bin at 0.4--2\,GeV/$c$.}
\caption{MuDst mixed event (main result), 0--80\%.}
\end{figure}

Random matches are \ovtFrCi\% of the TOF matches in 0--5\% and \ovtFrAll\% in 0--80\% ($0.2<p<3$\,GeV/$c$). This agrees with
the emulation of Sec.~\ref{sec:muemul} (\muFrProdCi\% and \muFrProdAll\%). With the common calibration the spike of the
emulation beside the pion peak is gone: the largest bin-by-bin ME/SE at $0.4<p<1.5$\,GeV/$c$ is \ovtVmaxCi{} in 0--5\%, and
the mixed event follows the shape of the same-event background at all $m^2$, including the valleys and the region between
the p and d peaks (read by eye, Fig.~\ref{fig:ovtrk_m2_cent}). The fractions under the PID peaks are listed per momentum bin in
Table~\ref{tab:mudstcmp}.

\paragraph{Effect of the partner tracks on the event's own matches.} The added tracks also compete for the hits of the event,
which the overlay makes denser than a real event. Of the QA tracks matched in the same event, \ovtSEkeep\% keep their own hit
in the overlays, \ovtSEswap\% are matched to a partner hit instead and \ovtSElose\% lose their match, mostly because a
partner track shares the cell. Without partner tracks the overlay matches \ovAnyMatch\% of the QA tracks (same event
\ovSEMatch\%); with them \ovtAnyMatch\%.

\caveat{Between the p and d peaks the mixed event is still \ovtHiMinCi--\ovtHiMaxCi{} times below the same event per momentum
bin at 0.4--2\,GeV/$c$ (0--5\%), with \ovtNHiMin--\ovtNHiMax{} same-event entries per bin. Below 0.4\,GeV/$c$ and above 2\,GeV/$c$ the
ratio is larger (Table~\ref{tab:mudstcmp}), and these bins are sparse. The overlay removes \ovtSElose\% of the event's own
matches, which slightly reduces the competition and is a possible bias of the same order. Light-nucleus candidates and
other real tracks in the window are not subtracted. The larger sample is needed to tell these apart.}

\subsubsection{Cross-checks: overlay without partner tracks, and the pool method}
\label{sec:mupool}
The same maker was run in three other setups on the same events:
\begin{itemize}
\item \textbf{Overlay, partner hits only:} random matches \ovFrCi\% (0--5\%) and \ovFrAll\% (0--80\%), SE/ME between the p and
  d peaks \ovHiCi. Without their tracks, partner hits of reconstructed particles are never discarded as shared, so the
  random matches are overestimated by a factor \ovHitsOnlyFactor{} (0--5\%).
\item \textbf{Pool} (the conventional mixed event: the event's own hits removed, its tracks matched to the partner hits only):
  \poolFrCi\% and \poolFrAll\%, SE/ME \poolHiCi. Every track is treated as if its own hit were missing, so the result is
  the random-hit probability for tracks without their own hit, not the random-match fraction. The overshoot reaches
  ME/SE = \poolVmaxCi{} in single valley bins (0--5\%).
\item \textbf{Pool with partner tracks:} \pooltFrCi\% and \pooltFrAll\%, SE/ME \pooltHiCi. The partner tracks remove
  \poolTrkRemoved\% of the pool matches (0--5\%): the hits left free are mostly those of unreconstructed particles.
\end{itemize}

\begin{figure}[H]
\centering
\figpanel{0.49}{ovpool_m2_cent}{Pool: event hits removed, partner hits only.}
\figpanel{0.49}{ovpooltrk_m2_cent}{Pool with partner tracks.}
\caption{Pool setups, 0--5\%; encodings as in Fig.~\ref{fig:ovtrk_m2_cent} (\fpath{scripts/make_ov.C}).}
\end{figure}

\choice{The pool method cannot give the random-match fraction without the fraction of tracks that have no own hit, which is
not measurable (the observed matches contain random ones), and it misses random hits that beat the track's own hit by ToT.
The overlay keeps the own hit and lets the production algorithm decide both cases; it is used for the result.}

\subsubsection{Comparison and current result}
\input{tables/mudst_cmp}

\decision{Current result: the overlay of partner hits and tracks with the production makers re-run (method 2). Random
matches are \ovtFrCi\% of the TOF matches in 0--5\% and \ovtFrAll\% in 0--80\% ($0.2<p<3$\,GeV/$c$), with the emulation
(\muFrProdCi\% and \muFrProdAll\%) as the cross-check. To be confirmed on the 98 restored MuDst files of 14 runs.}

The MuDst mixed event describes the same-event background between the p and d peaks far better than any picoDst mixed event:
SE/ME there is \ovtHiCi{} in 0--5\%, against \mtSbRatioAllVnine{} on picoDst. The late background is therefore made of random
matches to real, mostly unreconstructed slow particles, which have no time in the picoDst. Above 2\,GeV/$c$ and below
0.4\,GeV/$c$ the region has few same-event entries per bin ($N_{\rm HI}$ in Table~\ref{tab:mudstcmp}); these bins need the
larger sample (98 MuDst files restored).
